\section{Detailed progress-measure methodology}
\label{app:detailed_methods}

This appendix mirrors Section~5 of \citet{nanda2023progress} with TRM-specific implementation notes.

\subsection{Logit grid extraction}
For each pair $(a,b)$ in the modular-addition dataset, TRM receives input tokens $[a, b, \texttt{=}]$.
We run a single forward pass with ACT halted and read logits at position~2 (the ``$=$'' token).
Stacking over all $P^2$ pairs yields $L \in \mathbb{R}^{P \times P \times P}$ where $L_{a,b,c}$ is the logit for output $c$.

\subsection{Fourier basis on cyclic groups}
For prime modulus $P=113$, define $\omega_k = 2\pi k / P$.
Basis functions on index $x \in \{0,\ldots,P-1\}$:
\begin{align}
  \phi_k^{\cos}(x) &= \cos(\omega_k x), &
  \phi_k^{\sin}(x) &= \sin(\omega_k x).
\end{align}
For the logit map, the trigonometric identity $\cos(\omega_k(a+b-c))$ decomposes into products of 1D Fourier components along $a$, $b$, and $c$.

\subsection{Restricted (trigonometric) loss}
Given key frequencies $K = \{k_1,\ldots,k_5\}$, we project $L$ onto the span of $\cos(\omega_k(a+b))$ and $\sin(\omega_k(a+b))$ directions for $k \in K$, apply bias correction per \citet{nanda2023progress}, and compute cross-entropy against labels on train and test splits.

\subsection{Excluded loss}
For each $k \in K$ in descending sensitivity order, we subtract the rank-1 component associated with frequency $k$ from logits and measure test CE.
High excluded loss at step~2000 and low excluded loss at convergence indicate circuit formation.

\subsection{Embedding and unembedding Fourier norms}
$W_E \in \mathbb{R}^{P \times d}$: apply 1D DFT along the token axis (rows), sum $\ell_2$ energy per frequency pair $(k_{\cos}, k_{\sin})$.
$W_U \in \mathbb{R}^{P \times d}$: same along the vocabulary axis of the output head.
Key frequencies are those with largest excluded-loss sensitivity (not simply largest embedding norm).

\subsection{Gini coefficient}
Following Nanda Appendix~B, we compute the Gini coefficient of the vector of embedding Fourier norms to track sparsification during grokking.

\section{Reverse-engineering methodology}
\label{app:reverse_methods}

\subsection{$W_L$ neuron-logit map}
For SwiGLU layer $\ell$ with down-projection $W_{\mathrm{down}}^{(\ell)} \in \mathbb{R}^{d \times d_{\mathrm{ff}}}$ and unembedding $W_U \in \mathbb{R}^{P \times d}$:
\begin{equation}
  W_L^{(\ell)} = (W_{\mathrm{down}}^{(\ell)})^\top W_U^\top \in \mathbb{R}^{d_{\mathrm{ff}} \times P}.
\end{equation}
Each row of $W_L$ maps one MLP neuron to logits over $c$; we DFT each row along the $P$-axis.

\subsection{MLP neuron clustering}
For all $(a,b)$ pairs, record post-activation MLP outputs $h_n(a,b)$ at the ``$=$'' token.
Assign neuron $n$ to frequency $k$ maximizing $R^2(h_n, \cos(\omega_k(a+b)))$ or $\sin(\omega_k(a+b))$.

\subsection{Attention decomposition}
For layer~0, extract mean attention weights from query position~2 (``$=$'') to positions~0 ($a$) and~1 ($b$), averaged over all $(a,b)$ and batch.
Compare per-head routing strength to Nanda's Figure~4.

\subsection{Component patching}
We patch attention outputs and MLP outputs separately, measuring change in test CE relative to baseline, following the causal scrubbing spirit of \citet{nanda2023progress} Section~4.3.

\section{Latent probe methodology}
\label{app:latent_methods}

\subsection{ACT outer loop}
TRM uses adaptive computation time with maximum 8 steps.
At each ACT step $t$, we record $z_H$ at the ``$=$'' token and the predicted logit distribution.

\subsection{Fixed-point violation}
For test examples where step~0 prediction is correct, we flag violation if any later step changes the argmax prediction.

\subsection{Reasoning modes}
Following \citet{ren2026reasoning}: \textbf{trivial success} (correct at step~0), \textbf{nontrivial success} (wrong then correct), \textbf{nontrivial failure} (never correct), \textbf{oscillation} (alternating predictions).

\subsection{PCA trajectories}
For individual examples, stack $z_H^{(t)} \in \mathbb{R}^d$ across ACT steps and apply PCA for visualization.
